\section{机器的“护身符”}
\markboth{23. 机器的“护身符”}{23. 机器的“护身符”}

卡丹圆偶可以使转动变成往返直线运动，而且从 $x=2r\cos\theta$ 看出，这种往返直线运动是遵循余弦规律的。这一特点可应用于机械设计。

我们知道对于沿一线段 $AG$ 作往返直线运动的物体，在它到达行程的端点 $A$ 或 $G$ 时，速度要改变方向，因而产生撞击力。在端点的速度越大，撞击力越强。特别是当机器运转很快时，这种撞击力的影响尤其明显，可以造成机器的某些部件迅速损坏。所以在设计机器时，对于作往返直线运动的部件，尤其是高速运动的部件，常需设法使它在行程端点的速度尽量小些，最好在端点的速度为零。

对于由卡丹圆偶产生的往返直线运动，位移规律是
\[
  x = 2r\cos\omega t,
\]
这里 $t$ 是时间，$\omega = \dfrac{\theta}{t}$ 是动圆圆心 $C$ 绕定圆圆心 $O$ 旋转的角速度。现在我们来看看它是否能满足上述要求。

将上式对 $t$ 求导两次，顺次得到在时刻 $t$ 的速度 $v$ 和加速度 $a$ 分别是
\begin{align*}
  v &= x' = -2r\omega\sin\omega t, \\
  a &= v' = -2r\omega^2\cos\omega t.
\end{align*}

在行程的中点处 $x=0$，因而 $\cos\omega t=0$，$\sin\omega t=\pm 1$，这时 $a=0$，$|v|$ 最大。随着动点逐渐向行程端点移动，$\cos\omega t$ 的绝对值逐渐增大，$\sin\omega t$ 的绝对值逐渐减小，因而速度 $v$ 逐渐减慢。直到到达行程端点时，$x=\pm 2r$，$\cos\omega t=\pm 1$，因而 $\sin\omega t=0$，$v=0$。这时加速度最大，速度为零，因而使撞击力大大削弱。由此看来，这“$\cos$”还是机器的“护身符”呢。

在机械设计中有很多种获得按余弦规律作往返直线运动的方法，利用卡丹圆偶只是其中的一种方法。
